\documentclass[11pt]{article}

\usepackage[T1]{fontenc}
\usepackage{lmodern}
\usepackage[margin=1in]{geometry}
\usepackage{booktabs}
\usepackage{longtable}
\usepackage{array}
\usepackage{tabularx}
\newcolumntype{Y}{>{\raggedright\arraybackslash}X}

\title{EPSS-TPG Experiment Results}
\author{}
\date{}

\begin{document}
\maketitle

\section{What This Document Covers}

This document reports the test-set results from the current round of
EPSS-TPG experiments. The model and the input pipelines are explained
in two companion documents that share the same vocabulary and
notation: one walks through the dataset features, the other walks
through the multi-view TPG formula. This document focuses on the
numbers.

The experiments fall into five groups:

\begin{enumerate}
    \item A 19-run text-source ablation on the four refreshed
    Sec4AI4Aec-EPSS-Enhanced datasets: GPT, Gemma, Mistral and DeepSeek.
    Mistral is new in this round; GPT and Gemma were re-fetched against
    the new schema; the older Llama variant was dropped because its
    schema is incompatible. The ablation should be 20 runs, but one
    cell is unrunnable (DeepSeek does not ship the CVSS-templated
    summary column), so 19 runs actually completed.
    \item A 15-run text-source ablation on the three Megavul-branch
    datasets: Mega-GPT, Mega-Gemma and Mega-Mistral. Megavul is a
    commit-grounded vulnerability corpus with a much sparser positive
    class than the social-media datasets (about $1.7\%$ versus
    $15.5\%$ at the same EPSS binarisation threshold).
    \item A cross-reference of the model's highest-confidence test
    predictions against the live CISA KEV catalog
    (catalogVersion 2026.05.08, $1{,}590$ CVEs) and the live FIRST
    EPSS API (queried 2026-05-12).
    \item An NVD/KEV binary classification rerun on the older balanced
    dataset.
    \item Two NVD/KEV temporal-shift evaluations: train on 2020--2022,
    test on 2023--2024.
\end{enumerate}

The two text-source ablations and the KEV cross-reference are the main
story; the NVD/KEV runs are kept for context.

\section{The Social-Media Text-Source Ablation}

\subsection{Setup}

The question we are answering: does feeding the model an LLM-generated
summary of a CVE help, hurt, or do nothing, compared with just feeding
it the CVE description? Each of the four social-media datasets ships
the same CVE description plus three different LLM summaries:

\begin{itemize}
    \item one summary built from every source link the dataset has for
    that CVE;
    \item one summary built only from GitHub URLs;
    \item one summary written from the CVSS metrics in templated form.
\end{itemize}

We trained the model under five settings per dataset, so the design is
$4 \times 5 = 20$ runs.

\begin{table}[h]
\centering
\footnotesize
\begin{tabularx}{\textwidth}{@{}l Y Y@{}}
\toprule
Variant & Flags passed at training time & What text the model actually sees \\
\midrule
\textsc{D}      & \texttt{-{}-summary-source description} & The CVE description, nothing else. The baseline. \\[2pt]
\textsc{S\_all} & \texttt{-{}-summary-only-tpg} \newline \texttt{-{}-summary-source all\_sources}  & Only the all-sources summary. CVEs without that summary are dropped, so the train and test sets shrink. \\[2pt]
\textsc{S\_git} & \texttt{-{}-summary-only-tpg} \newline \texttt{-{}-summary-source github\_urls}  & Only the GitHub-URL summary. Most CVEs do not have one, so the test set drops to $339$--$360$ samples instead of $1{,}385$. \\[2pt]
\textsc{S\_cvss}& \texttt{-{}-summary-only-tpg} \newline \texttt{-{}-summary-source cvss\_metrics} & Only the templated CVSS-metrics summary. Available for GPT, Gemma and Mistral; the DeepSeek dataset does not ship this column. \\[2pt]
\textsc{ALL}    & \texttt{-{}-include-summary-in-tpg} \newline \texttt{-{}-summary-source combined} & The description plus all three summaries, joined with double newlines. The largest input by token count. \\
\bottomrule
\end{tabularx}
\caption{The five settings used per dataset.}
\end{table}

Every run uses the same training command, with these flags:
\texttt{-{}-backbone multiview}, \texttt{-{}-hybrid},
\texttt{-{}-label-mode soft}, \texttt{-{}-epochs 100} and
\texttt{-{}-no-epss-feature}.
The \texttt{-{}-no-epss-feature} flag matters: the soft target is the EPSS
score itself, so the EPSS score and percentile must be removed from the
tabular inputs. Otherwise the model would just learn to copy the input
back as the prediction. Input channel size stays at $781$ ($13$-D
node-type one-hot plus $768$-D SecBERT embedding); the tabular vector
is $55$-D without EPSS; the edge-type vocabulary is the base $13$ (no
security edges in this matrix).

\paragraph{Why $19$ runs and not $20$.}
The DeepSeek dataset has no CVSS-templated summary column. Zero records
out of $9{,}218$. So the DeepSeek \textsc{S\_cvss} cell cannot be trained
and is reported as ``--'' below. This is a missing-data fact about the
DeepSeek dataset, not a training failure. The other $19$ runs all
completed.

\subsection{Per-Run Numbers}

\small
\begin{longtable}{llrrrrrrr}
\caption{Test-set metrics for the $19$ completed social-media runs.
The missing DeepSeek \textsc{S\_cvss} cell is shown as ``--''.}\\
\toprule
Dataset & Variant & PR-AUC & ROC-AUC & F1 & Precision & Recall & Brier & $N_{\text{test}}$ \\
\midrule
\endfirsthead
\toprule
Dataset & Variant & PR-AUC & ROC-AUC & F1 & Precision & Recall & Brier & $N_{\text{test}}$ \\
\midrule
\endhead
\bottomrule
\endfoot

GPT      & D       & 0.8438 & 0.9438 & 0.8056 & 0.8113 & 0.8000 & 0.0491 & 1{,}385 \\
GPT      & S\_all  & 0.6826 & 0.8991 & 0.6407 & 0.5426 & 0.7821 & 0.1004 & 1{,}143 \\
GPT      & S\_git  & 0.7048 & 0.9374 & 0.7846 & 0.7083 & 0.8793 & 0.0691 &   339 \\
GPT      & S\_cvss & 0.5358 & 0.8584 & 0.5437 & 0.5529 & 0.5349 & 0.1019 & 1{,}385 \\
GPT      & ALL     & 0.8003 & 0.9442 & 0.7440 & 0.7739 & 0.7163 & 0.0588 & 1{,}385 \\

Gemma    & D       & 0.8356 & 0.9338 & 0.8159 & 0.8178 & 0.8140 & 0.0505 & 1{,}385 \\
Gemma    & S\_all  & 0.7332 & 0.9176 & 0.5134 & 0.8272 & 0.3722 & 0.0762 & 1{,}148 \\
Gemma    & S\_git  & 0.7186 & 0.8968 & 0.5684 & 0.7500 & 0.4576 & 0.0913 &   360 \\
Gemma    & S\_cvss & 0.5421 & 0.8654 & 0.5558 & 0.5680 & 0.5442 & 0.1021 & 1{,}385 \\
Gemma    & ALL     & 0.8341 & 0.9494 & 0.7669 & 0.8315 & 0.7116 & 0.0533 & 1{,}385 \\

Mistral  & D       & 0.8378 & 0.9448 & 0.7990 & 0.8227 & 0.7767 & 0.0499 & 1{,}385 \\
Mistral  & S\_all  & 0.6273 & 0.8887 & 0.4901 & 0.6966 & 0.3780 & 0.0838 & 1{,}099 \\
Mistral  & S\_git  & 0.6266 & 0.8828 & 0.5319 & 0.7143 & 0.4237 & 0.0990 &   359 \\
Mistral  & S\_cvss & 0.5510 & 0.8697 & 0.5383 & 0.5321 & 0.5446 & 0.1031 & 1{,}354 \\
Mistral  & ALL     & 0.8313 & 0.9469 & 0.7656 & 0.7882 & 0.7442 & 0.0549 & 1{,}385 \\

DeepSeek & D       & 0.8326 & 0.9374 & 0.8000 & 0.7739 & 0.8279 & 0.0538 & 1{,}385 \\
DeepSeek & S\_all  & 0.6764 & 0.8922 & 0.6743 & 0.6941 & 0.6556 & 0.0771 & 1{,}147 \\
DeepSeek & S\_git  & 0.6937 & 0.8580 & 0.6061 & 0.7500 & 0.5085 & 0.0844 &   360 \\
DeepSeek & S\_cvss & --     & --     & --     & --     & --     & --     & --      \\
DeepSeek & ALL     & 0.8112 & 0.9377 & 0.7549 & 0.7979 & 0.7163 & 0.0579 & 1{,}385 \\

\end{longtable}
\normalsize

\subsection{Means by Variant}

\begin{table}[h]
\centering
\small
\begin{tabular}{lcrr}
\toprule
Variant & $n$ & Mean PR-AUC & Mean change vs \textsc{D} \\
\midrule
\textsc{D}      & 4 & 0.8374 & --       \\
\textsc{S\_all} & 4 & 0.6799 & $-0.1576$ \\
\textsc{S\_git} & 4 & 0.6859 & $-0.1515$ \\
\textsc{S\_cvss}& 3 & 0.5430 & $-0.2944$ \\
\textsc{ALL}    & 4 & 0.8192 & $-0.0182$ \\
\bottomrule
\end{tabular}
\caption{Mean PR-AUC across the four social-media datasets.
\textsc{S\_cvss} is over three datasets only because DeepSeek does
not ship that summary column.}
\end{table}

A small caveat before reading the deltas. \textsc{S\_all} and
\textsc{S\_git} drop CVEs whose chosen summary is empty, so their test
sets are smaller. \textsc{S\_git} is evaluated on $339$--$360$ samples
per run instead of $1{,}385$; \textsc{S\_all} on $1{,}099$--$1{,}148$.
The PR-AUC numbers are valid for the subsets they were measured on,
but the deltas against \textsc{D} are not perfectly like-for-like.

\subsection{Per-Dataset PR-AUC}

\begin{table}[h]
\centering
\small
\begin{tabular}{lrrrrr}
\toprule
Dataset & \textsc{D} & \textsc{S\_all} & \textsc{S\_git} & \textsc{S\_cvss} & \textsc{ALL} \\
\midrule
GPT      & \textbf{0.8438} & 0.6826 & 0.7048 & 0.5358 & 0.8003 \\
Gemma    & \textbf{0.8356} & 0.7332 & 0.7186 & 0.5421 & 0.8341 \\
Mistral  & \textbf{0.8378} & 0.6273 & 0.6266 & 0.5510 & 0.8313 \\
DeepSeek & \textbf{0.8326} & 0.6764 & 0.6937 & --     & 0.8112 \\
\midrule
Mean     & 0.8374          & 0.6799 & 0.6859 & 0.5430 & 0.8192 \\
\bottomrule
\end{tabular}
\caption{PR-AUC by dataset and variant. Bold marks the best variant
for that dataset. The description-only baseline wins on every dataset.}
\end{table}

\subsection{What the Numbers Say}

\paragraph{Description alone is the strongest single text source.}
The description-only baseline wins on all four datasets. The four
\textsc{D} runs sit in a band $0.011$ wide ($0.8326$ to $0.8438$).
That tightness is what we expected: the description column is
identical across all four datasets, only the summary columns change
with the LLM provider, so the four \textsc{D} runs are essentially
the same model trained four times with the same input.

\paragraph{LLM summaries on their own are worse than the description.}
\textsc{S\_all} loses about $0.16$ PR-AUC on average. \textsc{S\_git}
loses about $0.15$. \textsc{S\_cvss} is the worst at $-0.29$. Three
things are happening at once. First, the summaries are paraphrases:
they turn dense vulnerability text into looser narrative text. Second,
two of the three summaries are missing for many CVEs, so those runs
train on smaller datasets. Third, the CVSS-metrics summary is
templated CVSS prose like ``Severity: High; Attack Vector: NETWORK''.
The model already sees the same CVSS data as $22$ one-hot tabular
features, so reading it again as text adds no information and
appears to push the description signal out.

\paragraph{Adding all three summaries on top of the description is a wash.}
The \textsc{ALL} variant concatenates the description and the three
summaries. The mean PR-AUC is $0.8192$, just $0.018$ below the
description-only baseline. Per-dataset: GPT loses $0.044$, Gemma
loses $0.002$, Mistral loses $0.007$, DeepSeek loses $0.021$. The
most likely cause is that the extra summary tokens get mean-and-max
pooled into the same graph embedding as the description tokens, so
the description signal is diluted rather than reinforced.

\paragraph{The variant ordering is the same across providers.}
On every dataset the order is roughly
$\textsc{D} \approx \textsc{ALL} \gg \textsc{S\_all} \approx \textsc{S\_git} > \textsc{S\_cvss}$.
\textsc{S\_all} and \textsc{S\_git} swap on GPT and Gemma but the
broad shape holds. The four LLM providers agree on the relative
quality of their own summaries.

\paragraph{Calibration follows the same order.}
The Brier scores tell the same story as PR-AUC. \textsc{D} and
\textsc{ALL} sit in the $0.049$ to $0.059$ band. \textsc{S\_all} is
around $0.07$ to $0.10$. \textsc{S\_cvss} is the worst at $0.10$ to
$0.11$.

\subsection{One-Paragraph Summary}

After removing the EPSS score from the tabular inputs (the source of
earlier $0.99$ PR-AUC numbers, which were target leakage), the
description-only baseline reaches mean PR-AUC $0.8374$ on the four
refreshed social-media datasets. Replacing the description with any
single LLM summary drops PR-AUC by between $0.15$ and $0.29$. The
templated CVSS summary is the worst because it duplicates information
already in the CVSS tabular features. Adding all three summaries on
top of the description lands within $0.02$ PR-AUC of the
description-only baseline. The CVE description is the only text source
that materially helps; the LLM summaries do not add net value on this
in-distribution corpus.

These are in-distribution results on the data shipped with the
Sec4AI4Aec-EPSS-Enhanced repository. They are not deployment numbers.
An earlier cross-distribution check on a held-out slice of NVD CVEs
saw about an $11\times$ drop in PR-AUC, which is a separate problem
caused by sample-selection bias.

\section{The Megavul Text-Source Ablation}

\subsection{Setup}

The Megavul side runs the same five-way variant scheme as the
social-media side, but on the three Megavul-branch datasets:
Mega-GPT, Mega-Gemma and Mega-Mistral. Each dataset covers $9{,}486$
rows and $9{,}240$ unique CVEs. Some CVEs appear in several rows
because Megavul links each CVE to one or more vulnerable commits;
the pipeline deduplicates by CVE-ID when building the labelled
records.

The schema differs from the social-media side. There is no social-media
post column, no observation date, no source-platform column. Instead,
every row carries a commit URL, the pre-patch and post-patch function
code, and CVSS information in two flavours (the original NVD CVSS plus
an LLM-augmented v3 estimate when the original is v2-only). The three
LLM summaries are different too: a summary of the commit URL, a summary
of the vulnerable function code, and a summary of the CVSS metrics.
The five variants are:

\begin{table}[h]
\centering
\footnotesize
\begin{tabularx}{\textwidth}{@{}l Y Y@{}}
\toprule
Variant & Flags passed at training time & What text the model actually sees \\
\midrule
\textsc{D}        & \texttt{-{}-summary-source description} & The CVE description, nothing else. Baseline. \\[2pt]
\textsc{S\_url}   & \texttt{-{}-summary-only-tpg} \newline \texttt{-{}-summary-source commit\_url} & Only the commit-URL summary. \\[2pt]
\textsc{S\_code}  & \texttt{-{}-summary-only-tpg} \newline \texttt{-{}-summary-source code} & Only the LLM summary of the vulnerable function code. \\[2pt]
\textsc{S\_cvss}  & \texttt{-{}-summary-only-tpg} \newline \texttt{-{}-summary-source cvss\_metrics} & Only the CVSS-metrics summary. \\[2pt]
\textsc{ALL}      & \texttt{-{}-include-summary-in-tpg} \newline \texttt{-{}-summary-source combined} & The description plus all three summaries, joined with double newlines. \\
\bottomrule
\end{tabularx}
\caption{Megavul five-way variants. Same training command as the
social-media side. The soft target is the EPSS score, sourced from
the most-recent EPSS recalculation per CVE with fallback to the
publication-time EPSS.}
\end{table}

\paragraph{Why the prevalence is so much lower than on the social-media side.}
The Megavul EPSS distribution is heavily skewed. At the same EPSS
binarisation threshold of $0.10$ used on the social-media side, only
about $1.7\%$ of Megavul records are positive, against $\sim 15.5\%$
on the social-media side. The Megavul model therefore trains and is
tested on a roughly $7\times$ sparser positive class. The
random-baseline PR-AUC equals the prevalence, which is $\approx 0.022$
on the Megavul test sets, so any Megavul PR-AUC number should be read
as a multiple of $0.022$ rather than as an absolute score.

\subsection{Per-Run Numbers}

\small
\begin{longtable}{llrrrrrrr}
\caption{Test-set metrics for the $15$ completed Megavul runs.}\\
\toprule
Dataset & Variant & PR-AUC & ROC-AUC & F1 & Precision & Recall & Brier & $N_{\text{test}}$ \\
\midrule
\endfirsthead
\toprule
Dataset & Variant & PR-AUC & ROC-AUC & F1 & Precision & Recall & Brier & $N_{\text{test}}$ \\
\midrule
\endhead
\bottomrule
\endfoot

Mega-GPT     & D       & 0.3597 & 0.9232 & 0.4138 & 0.3243 & 0.5714 & 0.0717 & 974 \\
Mega-GPT     & S\_url  & 0.2136 & 0.9119 & 0.2708 & 0.1733 & 0.6190 & 0.0926 & 972 \\
Mega-GPT     & S\_code & 0.1722 & 0.8434 & 0.2222 & 0.2000 & 0.2500 & 0.1250 & 947 \\
Mega-GPT     & S\_cvss & 0.0700 & 0.7840 & 0.0000 & 0.0000 & 0.0000 & 0.0754 & 974 \\
Mega-GPT     & ALL     & 0.3723 & 0.9044 & 0.3860 & 0.3056 & 0.5238 & 0.0991 & 974 \\

Mega-Gemma   & D       & 0.4102 & 0.8984 & 0.3636 & 0.2941 & 0.4762 & 0.0605 & 974 \\
Mega-Gemma   & S\_url  & 0.1939 & 0.8823 & 0.3415 & 0.3500 & 0.3333 & 0.0717 & 974 \\
Mega-Gemma   & S\_code & 0.1756 & 0.8642 & 0.2308 & 0.1579 & 0.4286 & 0.1047 & 974 \\
Mega-Gemma   & S\_cvss & 0.0600 & 0.7689 & 0.1143 & 0.0816 & 0.1905 & 0.0826 & 974 \\
Mega-Gemma   & ALL     & 0.4617 & 0.9039 & 0.4571 & 0.5714 & 0.3810 & 0.0326 & 974 \\

Mega-Mistral & D       & 0.4978 & 0.9194 & 0.3448 & 0.2273 & 0.7143 & 0.0931 & 974 \\
Mega-Mistral & S\_url  & 0.2482 & 0.9159 & 0.3396 & 0.2812 & 0.4286 & 0.1048 & 974 \\
Mega-Mistral & S\_code & 0.1885 & 0.8633 & 0.1875 & 0.2727 & 0.1429 & 0.0707 & 974 \\
Mega-Mistral & S\_cvss & 0.0905 & 0.8633 & 0.1778 & 0.1538 & 0.2105 & 0.1015 & 956 \\
Mega-Mistral & ALL     & 0.3163 & 0.9289 & 0.2727 & 0.2609 & 0.2857 & 0.0643 & 974 \\

\end{longtable}
\normalsize

\subsection{Means by Variant}

\begin{table}[h]
\centering
\small
\begin{tabular}{lcrrr}
\toprule
Variant & $n$ & Mean PR-AUC & Lift over chance & Change vs \textsc{D} \\
\midrule
\textsc{D}      & 3 & 0.4226 & $19.6\times$ & --       \\
\textsc{S\_url} & 3 & 0.2186 & $10.1\times$ & $-0.2040$ \\
\textsc{S\_code}& 3 & 0.1787 & $\phantom{0}8.3\times$ & $-0.2439$ \\
\textsc{S\_cvss}& 3 & 0.0735 & $\phantom{0}3.4\times$ & $-0.3491$ \\
\textsc{ALL}    & 3 & 0.3834 & $17.8\times$ & $-0.0392$ \\
\bottomrule
\end{tabular}
\caption{Mean PR-AUC across the three Megavul datasets. Lift is the
ratio of mean PR-AUC to the random-baseline PR-AUC, which on Megavul
equals the test-set prevalence ($\approx 0.022$).}
\end{table}

\subsection{Per-Dataset PR-AUC}

\begin{table}[h]
\centering
\small
\begin{tabular}{lrrrrr}
\toprule
Dataset & \textsc{D} & \textsc{S\_url} & \textsc{S\_code} & \textsc{S\_cvss} & \textsc{ALL} \\
\midrule
Mega-GPT     & 0.3597 & 0.2136 & 0.1722 & 0.0700 & \textbf{0.3723} \\
Mega-Gemma   & 0.4102 & 0.1939 & 0.1756 & 0.0600 & \textbf{0.4617} \\
Mega-Mistral & \textbf{0.4978} & 0.2482 & 0.1885 & 0.0905 & 0.3163 \\
\midrule
Mean         & 0.4226 & 0.2186 & 0.1787 & 0.0735 & 0.3834 \\
\bottomrule
\end{tabular}
\caption{Megavul PR-AUC by dataset and variant. Bold marks the best
variant for that dataset. \textsc{D} or \textsc{ALL} wins on every
Megavul dataset, identical to the social-media pattern.}
\end{table}

\subsection{What the Numbers Say (Megavul)}

\paragraph{Same variant ordering as on the social-media side.}
On every Megavul dataset the order is
$\textsc{D} \approx \textsc{ALL} \gg \textsc{S\_url} > \textsc{S\_code} > \textsc{S\_cvss}$.
Mega-GPT and Mega-Gemma tilt to \textsc{ALL}; Mega-Mistral tilts to
\textsc{D}. The relative ordering matches what we saw on the
social-media side: the description is the only single text source
that materially helps, the templated CVSS summary is the worst, and
concatenating everything is within a few hundredths of the
description-only baseline. So the variant ordering is not an artefact
of the social-media corpus; it reproduces on a commit-grounded
vulnerability corpus too.

\paragraph{Lower absolute PR-AUC, higher lift over chance.}
Megavul PR-AUC numbers run from $0.07$ to $0.50$, much lower than the
$0.54$ to $0.84$ on the social-media side. This is mostly the
prevalence, not the model. The Megavul test set has about $21$
positives in $974$ samples (prevalence $0.022$), so the random-baseline
PR-AUC is $0.022$. The best Megavul run, Mega-Mistral \textsc{D} at
PR-AUC $0.498$, is about $23\times$ the chance baseline. The best
social-media run, GPT \textsc{D} at PR-AUC $0.844$, is about $5.4\times$
the chance baseline of $0.155$. In lift terms the Megavul model is
the stronger ranker; in absolute terms it operates in a much harder
data regime.

\paragraph{ROC-AUC stays high across all Megavul variants.}
Even the worst Megavul variant, Mega-Gemma \textsc{S\_cvss}, has
ROC-AUC $0.769$. The three best variants per LLM all clear $0.90$.
ROC-AUC is the better headline metric on data this sparse: it is
invariant to prevalence and answers the question that matters
operationally, ``does the model rank a true positive above a random
negative?''. Brier scores tell a similar story: the best Megavul
runs sit in the $0.03$ to $0.07$ band, comparable to the best
social-media runs.

\paragraph{F1 at the default $0.5$ threshold understates Megavul performance.}
The pipeline records F1 at the default decision threshold of $0.5$.
For Megavul, where positives are $\sim 2\%$, the predicted
probabilities cluster well below $0.5$, so $0.5$-threshold F1 reads
low ($0.18$ to $0.46$). When the threshold is tuned on the validation
split, which is the standard choice for imbalanced binary reporting,
the picture improves. The best Megavul run reaches F1 $\approx 0.45$
at a val-tuned threshold; the weakest still reaches F1 $> 0.06$. FIRST
itself ships EPSS with operational thresholds at $0.10$ and $0.20$,
not $0.50$, so $0.5$-threshold F1 is the wrong choice on these data
in the first place.

\subsection{One-Paragraph Summary (Megavul)}

The five-way text-source ablation reproduces on the Megavul
commit-grounded corpus. Description-only is the strongest single
source (mean PR-AUC $0.4226$, $19.6\times$ random-baseline lift); the
templated CVSS summary is the worst (mean $0.0735$, $3.4\times$ lift);
concatenating description plus all three summaries lands within
$0.04$ PR-AUC of description-only. ROC-AUC stays in $0.77$ to $0.93$
across the $15$ runs. The lower absolute PR-AUC against the
social-media side comes from the $7\times$ sparser positive class
($1.7\%$ versus $15.5\%$), not from a weaker ranker: in lift terms
Megavul is comparable or stronger.

\section{Top-10 Predictions vs the CISA KEV Catalog}

For each of the seven dataset/LLM pairs (four social-media, three
Megavul) we picked the best-PR-AUC variant for that LLM and pulled
the model's top-10 highest-confidence test predictions. Each predicted
CVE was then checked against the live CISA KEV catalog
(catalogVersion $2026.05.08$, released 2026-05-08, $1{,}590$ CVEs)
and against the live FIRST EPSS API (queried 2026-05-12).

The hit rate against KEV is much higher on the social-media side.
That is expected. The social-media datasets cover CVEs from
2020--2024; the KEV catalog itself was created in November $2021$;
KEV-listed CVEs are exactly the kind of widely discussed vulnerabilities
that produce the social-media activity these datasets index. Megavul
on the other hand is dominated by long-tail OpenSSL fixes from
2013--2016 that pre-date the KEV catalog by half a decade.

\paragraph{Stored EPSS versus live EPSS.}
The seven top-10 tables below show a column called Stored EPSS. This
is the EPSS value that was held in the dataset at the time the
labelled records were built, i.e.\ the actual soft target the model
was trained against. It is not the same as today's live FIRST EPSS
value. FIRST recalculates EPSS daily, so for older CVEs (especially
the OpenSSL bugs that dominate the Megavul tables) the live value
today can be substantially different from what the dataset held. The
two values are reconciled in Table~\ref{tab:non-kev-status}, which
queries the live FIRST API directly. We use the stored value in the
top-10 tables because that is the right reference for evaluating the
model's predictions: the model cannot be expected to anticipate EPSS
values that FIRST had not yet recomputed at training time.

\subsection{Social-Media Side (Description-Only Variant)}

\small
\begin{longtable}{llp{0.30\textwidth}rrl}
\caption{Top-10 highest-predicted CVEs per social-media LLM, all
using the description-only variant (best PR-AUC per LLM). The
predicted column is the model probability. The Stored EPSS column is
the EPSS value as it was held in the dataset at label-generation
time, i.e.\ the soft target the model was trained to predict; this
is not the same as the live FIRST EPSS shown later in
Table~\ref{tab:non-kev-status}, which can differ for older CVEs
because EPSS is recalculated daily.}\\
\toprule
LLM & Rank & CVE / Vendor / Product & Pred. & Stored EPSS & In KEV \\
\midrule
\endfirsthead
\toprule
LLM & Rank & CVE / Vendor / Product & Pred. & Stored EPSS & In KEV \\
\midrule
\endhead
\bottomrule
\endfoot

GPT      &  1 & CVE-2024-9474 (Palo Alto Networks PAN-OS)               & 0.9935 & 0.9748 & yes \\
GPT      &  2 & CVE-2024-23897 (Jenkins CLI)                            & 0.9934 & 0.9447 & yes \\
GPT      &  3 & CVE-2024-4040 (CrushFTP)                                & 0.9926 & 0.9427 & yes \\
GPT      &  4 & CVE-2024-4040 (CrushFTP)$^{\dagger}$                    & 0.9926 & 0.9427 & yes \\
GPT      &  5 & CVE-2024-4040 (CrushFTP)$^{\dagger}$                    & 0.9926 & 0.9441 & yes \\
GPT      &  6 & CVE-2023-3519 (Citrix NetScaler ADC / Gateway)          & 0.9924 & 0.9628 & yes \\
GPT      &  7 & CVE-2024-0204 (Fortra GoAnywhere MFT)                   & 0.9915 & 0.9332 & no  \\
GPT      &  8 & CVE-2024-0204 (Fortra GoAnywhere MFT)$^{\dagger}$       & 0.9915 & 0.9332 & no  \\
GPT      &  9 & CVE-2018-13379 (Fortinet FortiOS)                       & 0.9914 & 0.9693 & yes \\
GPT      & 10 & CVE-2023-4966 (Citrix NetScaler ADC / Gateway)          & 0.9913 & 0.9438 & yes \\
\midrule
Gemma    &  1 & CVE-2024-23897 (Jenkins CLI)                            & 0.9940 & 0.9447 & yes \\
Gemma    &  2 & CVE-2024-0204 (Fortra GoAnywhere MFT)                   & 0.9926 & 0.9332 & no  \\
Gemma    &  3 & CVE-2024-0204 (Fortra GoAnywhere MFT)$^{\dagger}$       & 0.9926 & 0.9332 & no  \\
Gemma    &  4 & CVE-2018-13379 (Fortinet FortiOS)                       & 0.9885 & 0.9693 & yes \\
Gemma    &  5 & CVE-2023-35082 (Ivanti EPMM / MobileIron Core)          & 0.9885 & 0.9447 & yes \\
Gemma    &  6 & CVE-2023-35082 (Ivanti EPMM / MobileIron Core)$^{\dagger}$ & 0.9885 & 0.9447 & yes \\
Gemma    &  7 & CVE-2023-30258 (MagnusBilling)                          & 0.9880 & 0.9256 & no  \\
Gemma    &  8 & CVE-2023-3519 (Citrix NetScaler ADC / Gateway)          & 0.9880 & 0.9628 & yes \\
Gemma    &  9 & CVE-2023-4966 (Citrix NetScaler ADC / Gateway)          & 0.9880 & 0.9438 & yes \\
Gemma    & 10 & CVE-2024-28987 (SolarWinds Web Help Desk)               & 0.9879 & 0.9589 & yes \\
\midrule
Mistral  &  1 & CVE-2024-23897 (Jenkins CLI)                            & 0.9974 & 0.9447 & yes \\
Mistral  &  2 & CVE-2024-9474 (Palo Alto Networks PAN-OS)               & 0.9947 & 0.9748 & yes \\
Mistral  &  3 & CVE-2023-3519 (Citrix NetScaler ADC / Gateway)          & 0.9933 & 0.9628 & yes \\
Mistral  &  4 & CVE-2024-28987 (SolarWinds Web Help Desk)               & 0.9884 & 0.9589 & yes \\
Mistral  &  5 & CVE-2009-3960 (Adobe BlazeDS)                           & 0.9881 & 0.9382 & yes \\
Mistral  &  6 & CVE-2023-50164 (Apache Struts 2)                        & 0.9875 & 0.9374 & no  \\
Mistral  &  7 & CVE-2018-13379 (Fortinet FortiOS)                       & 0.9869 & 0.9693 & yes \\
Mistral  &  8 & CVE-2023-34362 (Progress MOVEit Transfer)               & 0.9854 & 0.9448 & yes \\
Mistral  &  9 & CVE-2023-49103 (ownCloud graphapi)                      & 0.9847 & 0.9437 & yes \\
Mistral  & 10 & CVE-2023-49103 (ownCloud graphapi)$^{\dagger}$          & 0.9847 & 0.9089 & yes \\
\midrule
DeepSeek &  1 & CVE-2024-23897 (Jenkins CLI)                            & 0.9965 & 0.9447 & yes \\
DeepSeek &  2 & CVE-2023-34039 (VMware Aria Operations for Networks)    & 0.9940 & 0.9325 & no  \\
DeepSeek &  3 & CVE-2024-8963 (Ivanti Cloud Services Appliance)         & 0.9929 & 0.9680 & yes \\
DeepSeek &  4 & CVE-2024-4040 (CrushFTP)                                & 0.9920 & 0.9441 & yes \\
DeepSeek &  5 & CVE-2024-4040 (CrushFTP)$^{\dagger}$                    & 0.9920 & 0.9427 & yes \\
DeepSeek &  6 & CVE-2024-4040 (CrushFTP)$^{\dagger}$                    & 0.9920 & 0.9427 & yes \\
DeepSeek &  7 & CVE-2023-3519 (Citrix NetScaler ADC / Gateway)          & 0.9913 & 0.9628 & yes \\
DeepSeek &  8 & CVE-2023-46805 (Ivanti Connect Secure / Policy Secure)  & 0.9908 & 0.9440 & yes \\
DeepSeek &  9 & CVE-2024-9474 (Palo Alto Networks PAN-OS)               & 0.9907 & 0.9748 & yes \\
DeepSeek & 10 & CVE-2022-3236 (Sophos Firewall)                         & 0.9906 & 0.9286 & yes \\

\end{longtable}
\normalsize

Rows marked with $^{\dagger}$ are repeat appearances of the same CVE
in the test set. The source datasets carry multiple rows per CVE
because the same CVE was discussed by several different sources
(Telegram, Reddit, Hacker News and so on), and each source-mention
becomes its own test sample. The model produces nearly identical
predictions for them, and they cluster at the top of the ranking
because they share the same CVE description text. We show all $10$
ranks rather than de-duplicating, so the table reflects exactly what
the model output.

Counting unique CVE-IDs, each LLM's top-10 contains $7$ to $9$
distinct CVEs and matches KEV on $6$ to $9$ of them. The four
distinct non-KEV CVEs across all four LLMs are: CVE-2024-0204
(Fortra GoAnywhere MFT, authentication bypass),
CVE-2023-30258 (MagnusBilling, unauthenticated command injection),
CVE-2023-50164 (Apache Struts 2, path traversal to RCE) and
CVE-2023-34039 (VMware Aria Operations for Networks, static SSH key
authentication bypass). All four have public exploit code or PoCs,
all four were the subject of multiple vendor advisories at disclosure
time, and all four currently sit in the top $0.3\%$ of FIRST EPSS, as
documented in Table~\ref{tab:non-kev-status}.

\subsection{Megavul Side (Best Variant Per LLM)}

\small
\begin{longtable}{llp{0.30\textwidth}rrl}
\caption{Top-10 highest-predicted CVEs per Megavul LLM, using the
best-PR-AUC variant for that LLM (\textsc{ALL} for Mega-GPT and
Mega-Gemma, \textsc{D} for Mega-Mistral). The Stored EPSS column is
the dataset-stored value, the same EPSS snapshot the model trained
against. For the older OpenSSL CVEs dominating this list, the live
FIRST EPSS today can be substantially lower or higher; see
Table~\ref{tab:non-kev-status} for the live values.}\\
\toprule
LLM & Rank & CVE / Vendor / Product & Pred. & Stored EPSS & In KEV \\
\midrule
\endfirsthead
\toprule
LLM & Rank & CVE / Vendor / Product & Pred. & Stored EPSS & In KEV \\
\midrule
\endhead
\bottomrule
\endfoot

Mega-Gemma   &  1 & CVE-2014-3505 (OpenSSL DTLS double free)            & 0.8090 & 0.9257 & no \\
Mega-Gemma   &  2 & CVE-2014-0195 (OpenSSL DTLS heap overflow)          & 0.7944 & 0.9676 & no \\
Mega-Gemma   &  3 & CVE-2014-0160 (OpenSSL Heartbleed)                  & 0.6915 & 0.9759 & yes \\
Mega-Gemma   &  4 & CVE-2016-6302 (OpenSSL ticket-length checks)        & 0.6712 & 0.1520 & no \\
Mega-Gemma   &  5 & CVE-2014-3470 (OpenSSL ECDH null deref)             & 0.6591 & 0.9737 & no \\
Mega-Gemma   &  6 & CVE-2014-3515 (OpenSSL ASN.1 type confusion)        & 0.6347 & 0.9433 & no \\
Mega-Gemma   &  7 & CVE-2013-0166 (OpenSSL handshake DoS)               & 0.6258 & 0.0068 & no \\
Mega-Gemma   &  8 & CVE-2015-0290 (OpenSSL multiblock corrupt)          & 0.5894 & 0.0866 & no \\
Mega-Gemma   &  9 & CVE-2015-0289 (OpenSSL PKCS7 null deref)            & 0.5771 & 0.0220 & no \\
Mega-Gemma   & 10 & CVE-2014-3569 (OpenSSL no-ssl3 build null deref)    & 0.5727 & 0.9025 & no \\
\midrule
Mega-GPT     &  1 & CVE-2014-3505 (OpenSSL DTLS double free)            & 0.8889 & 0.9257 & no \\
Mega-GPT     &  2 & CVE-2015-1792 (OpenSSL CMS null deref)              & 0.8384 & 0.0756 & no \\
Mega-GPT     &  3 & CVE-2015-0290 (OpenSSL multiblock corrupt)          & 0.8154 & 0.0866 & no \\
Mega-GPT     &  4 & CVE-2014-0160 (OpenSSL Heartbleed)                  & 0.7858 & 0.9759 & yes \\
Mega-GPT     &  5 & CVE-2016-2181 (OpenSSL DTLS replay protection)      & 0.7786 & 0.3711 & no \\
Mega-GPT     &  6 & CVE-2014-0195 (OpenSSL DTLS heap overflow)          & 0.7393 & 0.9676 & no \\
Mega-GPT     &  7 & CVE-2016-6302 (OpenSSL ticket-length checks)        & 0.7139 & 0.1520 & no \\
Mega-GPT     &  8 & CVE-2014-3569 (OpenSSL no-ssl3 build null deref)    & 0.7062 & 0.9025 & no \\
Mega-GPT     &  9 & CVE-2014-3470 (OpenSSL ECDH null deref)             & 0.6961 & 0.9737 & no \\
Mega-GPT     & 10 & CVE-2015-0289 (OpenSSL PKCS7 null deref)            & 0.6892 & 0.0220 & no \\
\midrule
Mega-Mistral &  1 & CVE-2014-3505 (OpenSSL DTLS double free)            & 0.9448 & 0.9257 & no \\
Mega-Mistral &  2 & CVE-2014-0195 (OpenSSL DTLS heap overflow)          & 0.9256 & 0.9676 & no \\
Mega-Mistral &  3 & CVE-2014-0160 (OpenSSL Heartbleed)                  & 0.9029 & 0.9759 & yes \\
Mega-Mistral &  4 & CVE-2016-2181 (OpenSSL DTLS replay protection)      & 0.8955 & 0.3711 & no \\
Mega-Mistral &  5 & CVE-2014-3470 (OpenSSL ECDH null deref)             & 0.8937 & 0.9737 & no \\
Mega-Mistral &  6 & CVE-2016-6302 (OpenSSL ticket-length checks)        & 0.8896 & 0.1520 & no \\
Mega-Mistral &  7 & CVE-2014-3569 (OpenSSL no-ssl3 build null deref)    & 0.8618 & 0.9025 & no \\
Mega-Mistral &  8 & CVE-2013-0166 (OpenSSL handshake DoS)               & 0.8412 & 0.0068 & no \\
Mega-Mistral &  9 & CVE-2014-2523 (Linux kernel netfilter DCCP)         & 0.8412 & 0.0747 & no \\
Mega-Mistral & 10 & CVE-2015-1790 (OpenSSL PKCS7 null deref)            & 0.8303 & 0.0756 & no \\

\end{longtable}
\normalsize

The Megavul top-10s overlap heavily across the three LLMs because the
Megavul corpus has a small set of OpenSSL bugs from 2013--2016 that
dominate the high-EPSS tail of the test set. Only one CVE
(CVE-2014-0160, Heartbleed) is in the current CISA KEV catalog. This
is a KEV-coverage issue, not a model issue: the KEV catalog started
accepting entries in November $2021$, while the Megavul OpenSSL bugs
are mostly from $2013$--$2016$ and so are simply outside the KEV
catalog's effective time window.

\subsection{Status of the Non-KEV Top-10 CVEs}
\label{tab:non-kev-status}

For each unique non-KEV CVE in any of the seven top-10 lists we
queried the live FIRST EPSS API and a brief web search for current
exploitation evidence. The $17$ unique non-KEV CVEs split into two
groups: four recent CVEs (2023--2024) where the model surfaces
high-confidence non-KEV predictions that are very plausibly being
exploited, and thirteen older OpenSSL or Linux-kernel CVEs
(2013--2016) that pre-date the KEV catalog.

\begin{table}[h]
\centering
\small
\begin{tabular}{lrrp{0.50\textwidth}}
\toprule
CVE & Live EPSS & Percentile & Exploitation evidence (May 2026) \\
\midrule
\multicolumn{4}{l}{\textit{Recent (2023--2024), not in KEV but high EPSS}} \\
\midrule
CVE-2024-0204  & 0.9305 & 99.79\% & Fortra GoAnywhere MFT auth bypass; PoC public January 2024; Shadowserver observed exploit attempts; advisories from Tenable, Rapid7, Help Net Security and others. \\
CVE-2023-50164 & 0.9286 & 99.77\% & Apache Struts 2 path traversal to RCE; PoC published December 2023; active exploitation reported by Akamai, Trend Micro, Cybersecurity Dive. \\
CVE-2023-34039 & 0.9317 & 99.80\% & VMware Aria Operations for Networks static SSH key, authentication bypass; PoC by Sina Kheirkhah; Broadcom CVSS 9.8 advisory. \\
CVE-2023-30258 & 0.9368 & 99.85\% & MagnusBilling unauthenticated command injection; CVSS 9.8; ranked at the $99.85$th EPSS percentile. \\
\midrule
\multicolumn{4}{l}{\textit{Older (2013--2016), pre-date KEV catalog}} \\
\midrule
CVE-2014-0195  & 0.9275 & 99.76\% & OpenSSL DTLS fragment heap overflow, RCE class. Multiple Metasploit modules historically. \\
CVE-2014-3470  & 0.9140 & 99.67\% & OpenSSL ECDH client null deref. \\
CVE-2014-3515  & 0.4866 & 97.78\% & OpenSSL ASN.1 type confusion. \\
CVE-2014-3505  & 0.4688 & 97.70\% & OpenSSL DTLS double free. \\
CVE-2015-0290  & 0.3256 & 96.90\% & OpenSSL multiblock corrupt. \\
CVE-2016-2181  & 0.2427 & 96.14\% & OpenSSL DTLS replay protection. \\
CVE-2015-1792  & 0.1236 & 93.95\% & OpenSSL CMS null deref. \\
CVE-2016-6302  & 0.1038 & 93.27\% & OpenSSL ticket-length checks. \\
CVE-2015-1790  & 0.0961 & 92.94\% & OpenSSL PKCS7 null deref. \\
CVE-2013-0166  & 0.0951 & 92.90\% & OpenSSL handshake DoS. \\
CVE-2014-3569  & 0.0756 & 91.89\% & OpenSSL no-ssl3 build null deref. \\
CVE-2015-0289  & 0.0579 & 90.57\% & OpenSSL PKCS7 null deref (related). \\
CVE-2014-2523  & 0.0463 & 89.36\% & Linux kernel netfilter DCCP. \\
\bottomrule
\end{tabular}
\caption{Live status of every non-KEV CVE that appears in any of the
seven top-10 prediction lists. Live EPSS and percentile from the FIRST
EPSS API on 2026-05-12.}
\end{table}

\paragraph{Reading.}
Of the $17$ non-KEV CVEs, all $17$ are in the top $11\%$ of FIRST EPSS
percentiles (the lowest is $89.36\%$). $12$ of $17$ have live EPSS
$\geq 0.10$, the same threshold the pipeline uses to binarise the
soft target. The four recent non-KEV CVEs all sit above the $99.7$th
percentile and have public exploit code or PoCs. The thirteen older
CVEs pre-date the November $2021$ start of the CISA KEV catalog. Put
together, the model's high-confidence non-KEV predictions are not
false positives. They are CVEs that FIRST itself ranks at the top of
its exploitation-likelihood scale; the model just identified them
through patterns in the description text.

\paragraph{Sources for the recent-CVE entries.}
Exploitation evidence for the four recent CVEs is aggregated from the
following public sources, all accessed 2026-05-12:
CVE-2024-0204 (Tenable, Rapid7, Help Net Security, UpGuard);
CVE-2023-50164 (Akamai, Trend Micro, Cybersecurity Dive, Palo Alto
Networks Unit 42, Qualys ThreatPROTECT);
CVE-2023-34039 (Help Net Security, Security Boulevard, Arctic Wolf,
The Hacker News);
CVE-2023-30258 (Wiz vulnerability database, CVE Details). Live EPSS
scores were retrieved from the FIRST EPSS API.

\section{NVD/KEV Binary Rerun}

This run is separate from the two text-source ablations. It uses the
older balanced NVD/KEV dataset. The target here is the binary KEV
membership label: $1$ means the CVE is in the CISA KEV catalog
snapshot used at label-generation time, $0$ means it is not.

The training command uses these flags: \texttt{-{}-backbone multiview},
\texttt{-{}-hybrid}, \texttt{-{}-label-mode binary} and
\texttt{-{}-skip-collect}. This run does not use any LLM summary and
does not enable security
edges. Unlike the EPSS soft-label runs, it keeps the EPSS score and
percentile in the tabular branch. That is not direct target leakage
(the label is KEV membership, not EPSS), but EPSS is a strong external
risk signal, so this should be read as a hybrid KEV predictor with
external risk input rather than as a text-only result.

\begin{table}[h]
\centering
\small
\begin{tabular}{ll}
\toprule
Property & Value \\
\midrule
Dataset family & older balanced NVD/KEV dataset \\
Dataset size & $4{,}015$ CVEs \\
Test samples & $604$ \\
Test positives & $121$ \\
Test prevalence & $20.03\%$ \\
Backbone & multiview \\
Model type & hybrid (graph + tabular) \\
Input channels & $781$ \\
Tabular dimensions & $57$ (EPSS kept as input here) \\
Edge types & $13$ \\
Epochs completed & $24$ \\
Best validation PR-AUC & $0.9247$ at epoch $9$ \\
Best validation F1 & $0.8511$ at epoch $10$ \\
\bottomrule
\end{tabular}
\caption{NVD/KEV rerun setup.}
\end{table}

\begin{table}[h]
\centering
\small
\begin{tabular}{lrrrrrrr}
\toprule
Run & PR-AUC & ROC-AUC & F1 & Precision & Recall & Brier & Loss \\
\midrule
Previous saved run & 0.7592 & 0.8923 & 0.6917 & 0.6975 & 0.6860 & 0.1073 & 2.0044 \\
Fresh rerun         & 0.9375 & 0.9749 & 0.8522 & 0.8991 & 0.8099 & 0.0479 & 0.9002 \\
Change              & $+0.1783$ & $+0.0827$ & $+0.1605$ & $+0.2016$ & $+0.1240$ & $-0.0594$ & $-1.1042$ \\
\bottomrule
\end{tabular}
\caption{Fresh NVD/KEV rerun versus the previous saved run.}
\end{table}

At threshold $0.5$ the test confusion matrix is
$\mathrm{TP} = 98$, $\mathrm{FP} = 11$, $\mathrm{TN} = 472$,
$\mathrm{FN} = 23$. The best-F1 threshold from the saved predictions
is around $0.068$. At that threshold F1 rises to about $0.8735$, with
precision $0.8629$ and recall $0.8843$. The table reports
$0.5$-threshold metrics because those are what the pipeline records
by default.

The fresh rerun is much stronger than the older saved result. The
most likely cause is that the rerun was built on a freshly processed
dataset with the current $57$-dim tabular vector, while the older run
used a stale graph cache that produced only $53$ tabular dimensions.

\section{NVD/KEV Temporal Shift}

This run measures temporal generalisation. The model is trained and
validated on CVEs published between $2020$ and $2022$, then tested on
a separate $2023$--$2024$ slice. This is harder than a random split
because the test CVEs come from years the model has not seen.

The target is still KEV membership. The model has access to CVE text,
CVSS and CWE metadata, EPSS fields, and exploit and source features.

\begin{table}[h]
\centering
\small
\begin{tabular}{ll}
\toprule
Property & Value \\
\midrule
Train window & 2020-01-02 to 2022-12-30 \\
Test window  & 2023-01-01 to 2024-12-30 \\
Train records & $7{,}487$ ($487$ positive, $6.50\%$) \\
Test records  & $3{,}316$ ($316$ positive, $9.53\%$) \\
Backbone & multiview \\
Model type & hybrid (graph + tabular) \\
Tabular dimensions & $57$ \\
Edge types & $13$ \\
EPSS in tabular input & yes \\
Security edges & no \\
Epochs completed & $27$ \\
Best validation PR-AUC & $0.8499$ at epoch $7$ \\
Best validation F1 & $0.7692$ at epoch $18$ \\
\bottomrule
\end{tabular}
\caption{Temporal-shift setup.}
\end{table}

\subsection{Same Window With EPSS Removed}

A follow-up run trained on the same $2020$--$2022$ window and tested
on the same $2023$--$2024$ window, but with the EPSS score and
percentile removed from the tabular inputs (\texttt{-{}-no-epss-feature}).
The setup is otherwise identical. This run measures how much of the
temporal-shift performance comes from the EPSS input.

\begin{table}[h]
\centering
\small
\begin{tabular}{lrrrrrr}
\toprule
Setting & PR-AUC & ROC-AUC & F1 & Precision & Recall & Brier \\
\midrule
EPSS in tabular   & 0.8456 & 0.9598 & 0.7806 & 0.8253 & 0.7405 & 0.0676 \\
EPSS removed      & 0.7041 & 0.9213 & 0.6519 & 0.7421 & 0.5808 & 0.0921 \\
Change            & $-0.1415$ & $-0.0385$ & $-0.1287$ & $-0.0832$ & $-0.1597$ & $+0.0245$ \\
\bottomrule
\end{tabular}
\caption{Temporal-shift run with EPSS in versus out of the tabular
branch. Same train/test split.}
\end{table}

Removing EPSS from the input drops PR-AUC by $0.14$ on the temporal
test set. The drop is real but the model still ranks well above
random (PR-AUC $0.70$ at prevalence $0.095$ is about $7.4\times$
chance). On this dataset the text and CVSS features carry most of the
signal; EPSS is a useful but not load-bearing additional input.

\section{NVD/KEV No-EPSS Temporal Extrapolation}

A third NVD/KEV run extrapolates further: train on $2020$ only, test
on $2021$--$2026$. EPSS is removed from the tabular input. This is
the hardest temporal evaluation: small training window, very wide
test window, no EPSS hint.

\begin{table}[h]
\centering
\small
\begin{tabular}{ll}
\toprule
Property & Value \\
\midrule
Train window & 2020-01-02 to 2020-12-30 \\
Test window  & 2021-01-04 to 2026-12-30 \\
Train records & $1{,}500$ \\
Test records  & $9{,}303$ \\
Test positives & $785$ ($8.44\%$) \\
EPSS in tabular input & no \\
\bottomrule
\end{tabular}
\caption{No-EPSS temporal-extrapolation setup.}
\end{table}

\begin{table}[h]
\centering
\small
\begin{tabular}{lrrrrrr}
\toprule
Run & PR-AUC & ROC-AUC & F1 & Precision & Recall & Brier \\
\midrule
2020 train, 2021--2026 test, no EPSS & 0.4853 & 0.8773 & 0.4860 & 0.6018 & 0.4076 & 0.0738 \\
\bottomrule
\end{tabular}
\caption{No-EPSS temporal extrapolation.}
\end{table}

PR-AUC drops to $0.49$ on this evaluation, but is still about
$5.8\times$ the chance baseline of $0.084$. ROC-AUC stays at $0.88$,
which means the ranking quality survives the temporal shift even
without EPSS as an input. The big change is in calibration and recall:
the model only catches $40.8\%$ of the test-set positives at the
default threshold. A lower decision threshold would recover most of
the recall at the cost of precision; the val-tuned threshold for this
run is $0.226$, well below $0.5$.

\end{document}
